\documentclass[11pt,a4paper]{article}
\usepackage[margin=2.2cm]{geometry}
\usepackage{amsmath}
\usepackage{booktabs}
\usepackage{graphicx}
\usepackage{xcolor}
\usepackage{hyperref}

% Red placeholder. Search for \todo before submitting.
\newcommand{\todo}[1]{\textcolor{red}{\textbf{[TODO: #1]}}}

\title{CS6886 Assignment 5: Optimising Attention on a CPU}
\author{Yukash Ram CS23B069 \\ \small \url{https://github.com/Yukash09/CS6886_Assignment_5}}
\date{}

\begin{document}
\maketitle

\noindent All measurements were taken with the laptop on AC power, using clang 21.1.8 with
\texttt{-O3 -march=native}. Each reported time is the driver's median of 5 runs after one
warm-up run, and GFLOP/s is the driver's $4N^2d/t$.

% ===========================================================================
\section*{Q1: Machine}

\begin{center}
\begin{tabular}{@{}ll@{}}
\toprule
CPU & AMD Ryzen 5 7520U, Zen 2 (Mendocino) \\
Cores & 4 physical, 8 logical (SMT), no P/E split \\
Clock & base 2.8\,GHz, max 4.3\,GHz, all-core under load 4.07\,GHz \\
SIMD & AVX2 + FMA3, 256-bit (8 FP32), 2 FMA units/core, 16 YMM registers \\
Caches & L1d 32\,KiB/core, L2 512\,KiB/core, L3 4\,MiB shared, 64\,B lines \\
Memory & 16\,GB LPDDR5-5500, 2 channels $\times$ 32-bit \\
OS & Fedora 43, kernel 7.2.13 \\
\bottomrule
\end{tabular}
\end{center}

\noindent Most fields come from \texttt{lscpu}, \texttt{/sys/devices/system/cpu} and
\texttt{dmidecode}; the base clock is from AMD's product page and the FMA count from
\todo{cite Zen 2 FMA source}. The all-core clock was read from \texttt{scaling\_cur\_freq}
while all 8 hardware threads were busy, and again under the AVX2 load of the OpenMP variant
(4.03--4.16\,GHz with 4 or 8 threads, 4.24\,GHz with one).

\paragraph{Peak FP32.} A 256-bit FMA operates on 8 floats and counts as 2 FLOP per float, and
each core has two FMA units, so a core can do 32 FLOP per cycle. Using the single-core boost
clock for one core and the measured all-core clock for four,
\[
  P_1 = 32 \times 4.3\,\text{GHz} = 137.6\ \text{GFLOP/s}, \qquad
  P_4 = 4 \times 32 \times 4.07\,\text{GHz} = 521\ \text{GFLOP/s}.
\]
Only physical cores are counted, since the two SMT threads of a core share its FMA units.

\paragraph{Memory bandwidth.} The memory bus is 64 bits wide (two 32-bit channels), so at
5500\,MT/s the theoretical bandwidth is $5500 \times 10^6 \times 8\,\text{B} = 44.0$\,GB/s.
STREAM Triad, run with 76\,MiB arrays so that nothing fits in the L3, measures
$B_1 = 8.06$\,GB/s with one thread and $B_4 = 15.5$\,GB/s with four. These measured values are
what the code can actually reach, so they are the ones used for the roofline.

\paragraph{Ridge point.} A kernel with arithmetic intensity $I$ (FLOP per byte of DRAM traffic)
can reach at most $\min(P,\ I \cdot B)$, and the two limits meet at $I^\ast = P/B$:
\[
  I^\ast_1 = 137.6 / 8.06 = 17.1\ \text{FLOP/B}, \qquad
  I^\ast_4 = 521 / 15.5 = 33.5\ \text{FLOP/B}.
\]

% ===========================================================================
\section*{Q2: Naive baseline}

\begin{center}
\begin{tabular}{@{}lrrr@{}}
\toprule
Input & Time (ms) & GFLOP/s & \% of $P_1$ \\
\midrule
$N=1024,\ d=64$  & 129.5  & 2.07 & 1.50 \\
$N=2048,\ d=64$  & 524.8  & 2.05 & 1.49 \\
$N=4096,\ d=64$  & 2163.0 & 1.99 & 1.45 \\
$N=4096,\ d=128$ & 5146.5 & 1.67 & 1.21 \\
\bottomrule
\end{tabular}
\end{center}

\paragraph{Scaling.} The work grows as $N^2d$, so doubling $N$ should cost $\times4$ and
doubling $d$ should cost $\times2$. Doubling $N$ costs $\times4.05$ and $\times4.12$, which means
the cost per FLOP barely changes with $N$: the inner loops are limited by the core (the
dependent FMA chain described below) rather than by memory. Doubling $d$ costs $\times2.38$, and
the two matmul steps slow down for different reasons. To separate the effect of $d$ from the
effect of data size, I also timed random inputs at $N=1024,\ d=128$. $QK^\top$ runs at about
2.0\,GFLOP/s at $d=128$ for both $N=1024$ and $N=4096$, so its slowdown does not depend on data
size; it comes from the longer chain of dependent FMAs in each dot product, which leaves less
overlap between consecutive dot products. $PV$, on the other hand, grows by $\times2.10$ from
$d=64$ to $128$ at $N=1024$ but by $\times2.53$ at $N=4096$, which points to memory. Hardware
counters confirm this: at $N=4096$, going to $d=128$ raises the L2 miss rate from 53\% to 74\%,
and both L2 misses and L1 TLB misses per FLOP roughly double. This follows from the larger
stride ($4d = 512$\,B) with which step 3 walks down the columns of $V$.

\paragraph{Where the time goes.} I put timers around the three steps in a copy of
\texttt{attn\_naive.c} (median of 5 runs):

\begin{center}
\begin{tabular}{@{}lrrr@{}}
\toprule
Input & $QK^\top$ (ms) & softmax (ms) & $PV$ (ms) \\
\midrule
$N=1024,\ d=64$  & 46.7   & 3.9  & 78.9   \\
$N=2048,\ d=64$  & 193.9  & 16.4 & 319.2  \\
$N=4096,\ d=64$  & 799.7  & 65.2 & 1306.8 \\
$N=4096,\ d=128$ & 2150.2 & 64.9 & 3216.4 \\
\bottomrule
\end{tabular}
\end{center}

\noindent Across all inputs, $PV$ takes about 60\% of the time, $QK^\top$ about 37\% and softmax
only 3\%. The main reason both matmul steps are slow is that they are not vectorised. clang
reports that it \emph{cannot prove it is safe to reorder floating-point operations} in the
\texttt{acc += \dots} loops, and the disassembly confirms that the hot loop is a scalar
\texttt{vfmadd231ss}. Each of these FMAs does 2 FLOP and has to wait about 5 cycles for the
previous one to finish, which limits the loop to roughly 0.4 FLOP per cycle against a peak of 32.
\texttt{perf stat} agrees, giving an IPC of 1.45 and about 0.49 FLOP per cycle at $N=1024$,
which is why naive reaches only about 1.5\% of peak.

$PV$ is slower than $QK^\top$ even though both do the same number of FLOPs, because of how they
access memory. $QK^\top$ reads $K$ along its rows with a 4-byte stride, so every byte of each
cache line it loads is used. $PV$ reads $V$ down its columns with a stride of $4d = 256$\,B, so
every load touches a new cache line and uses only 4 of its 64 bytes. Softmax is cheap despite
its scalar \texttt{expf} calls (clang reports \emph{call instruction cannot be vectorized})
because it does $O(N^2)$ work, compared with $O(N^2d)$ for each matmul. The $N \times N$ matrix
is not yet a bottleneck either: at $N=4096$, $S$ is 64\,MiB and makes about four passes through
DRAM, roughly 0.27\,GB or 33\,ms at 8.06\,GB/s, which is under 2\% of the runtime.

% ===========================================================================
\section*{Q3: Optimisations}
\subsection*{(a) Loop reordering}

Both matmuls in naive put the sum over the shared index innermost, which the compiler cannot
vectorise without reordering FP additions. In \texttt{attn\_loops.c} I moved that index out of
the inner loop, so that the inner loop updates a different output element on every iteration
and reads memory contiguously. For $PV$ the order becomes $i \to j \to k$, with the inner loop
\texttt{O[i][k] += S[i][j] * V[j][k]}. For $QK^\top$ no loop order is contiguous in $K$ while
avoiding a reduction, so I first copy $K$ into a transposed $d \times N$ array $K^\top$ (an
$O(Nd)$ step, 1.5\,ms at $N=4096$) and use $i \to k \to j$ with
\texttt{S[i][j] += Q[i][k] * Kt[k][j]}. clang now reports both inner loops as vectorised (width 8,
interleave 4). Since each output element still accumulates its terms in the same order as
naive, \texttt{make check} passes with an error of exactly 0.

\begin{center}
\begin{tabular}{@{}lrrrr@{}}
\toprule
Input & Time (ms) & GFLOP/s & \% of $P_1$ & Speedup \\
\midrule
$N=1024,\ d=64$  & 14.8  & 18.14 & 13.2 & $\times8.75$ \\
$N=2048,\ d=64$  & 64.6  & 16.62 & 12.1 & $\times8.12$ \\
$N=4096,\ d=64$  & 264.2 & 16.26 & 11.8 & $\times8.19$ \\
$N=4096,\ d=128$ & 420.8 & 20.41 & 14.8 & $\times12.23$ \\
\bottomrule
\end{tabular}
\end{center}

At $N=4096,\ d=64$ the steps now take 107\,ms ($QK^\top$, 20\,GFLOP/s), 81\,ms ($PV$,
27\,GFLOP/s) and 64\,ms (softmax), so softmax has grown from 3\% to 25\% of the time. The matmuls
are still far from peak for two reasons. First, every vector FMA in the inner loop needs two
loads and one store (the accumulator \texttt{S[i][j]} or \texttt{O[i][k]} lives in memory, not
in a register), and a Zen 2 core can do only two loads and one store per cycle, which caps the
loop at one FMA per cycle, half of peak, even when everything hits in L1. Second, each row $i$
streams the whole of $K^\top$ (or $V$), 1\,MiB at $N=4096,\ d=64$, which does not fit in the
512\,KiB L2, so it is re-read from L3 for every row. Register tiling and cache tiling in (b)
address these two limits.

\subsection*{(b) Cache and register tiling}

\texttt{attn\_loops.c} keeps every partial sum in memory, so each vector FMA needs two loads and
a store. In \texttt{attn\_tiled.c} each matmul instead computes an $8 \times 8$ block of the
output at a time in a local array \texttt{acc[8][8]}. Because the tile size is a compile-time
constant, clang unrolls the tile loops and keeps the block in eight vector registers; the inner
loop then loads one 8-wide row of $K^\top$ (or $V$) and broadcasts eight scalars of $Q$ (or $P$),
and the outputs are stored only once per tile. Eight independent accumulators also hide most of
the FMA latency. For cache tiling, rows are processed in blocks of 64, so that the small slice of
$K^\top$ or $V$ used by a tile is reused by all eight row tiles of the block while it is still
in L1. I measured blocks of 16 to 256 rows afterwards: 16 rows was fastest on every input, 2--8\%
faster than the 64 used here, and 256 the slowest. I also tried a
$4 \times 16$ tile, which was slightly faster at $N=4096,\ d=64$ (190\,ms against 199\,ms) but
much slower at $N=2048$ (63\,ms against 42\,ms), and padded the row stride of $K^\top$ to
$N+16$, which helped the $8 \times 8$ tile at $d=128$ (283\,ms against 338\,ms). Index
arithmetic has to be done in \texttt{size\_t}: with \texttt{int} indices the same code executed
$1.93\times$ more instructions and ran $1.7\times$ slower. \texttt{make check} passes with an
error of exactly 0, since every output element still sums its terms in the original order.

\begin{center}
\begin{tabular}{@{}lrrrrr@{}}
\toprule
Input & Time (ms) & GFLOP/s & \% of $P_1$ & vs.\ naive & vs.\ loops \\
\midrule
$N=1024,\ d=64$  & 11.5  & 23.38 & 17.0 & $\times11.28$ & $\times1.29$ \\
$N=2048,\ d=64$  & 41.7  & 25.72 & 18.7 & $\times12.57$ & $\times1.55$ \\
$N=4096,\ d=64$  & 226.8 & 18.94 & 13.8 & $\times9.54$  & $\times1.16$ \\
$N=4096,\ d=128$ & 331.8 & 25.89 & 18.8 & $\times15.51$ & $\times1.27$ \\
\bottomrule
\end{tabular}
\end{center}

The matmuls themselves are now much faster: at $N=1024$ they run at 49\,GFLOP/s ($QK^\top$) and
57\,GFLOP/s ($PV$), around 40\% of peak. The overall gain is smaller because the rest of the
program did not change. At $N=4096,\ d=64$ the steps take 72\,ms ($QK^\top$), 70\,ms (softmax)
and 42\,ms ($PV$), so softmax, with its scalar \texttt{expf}, is now 38\% of the time. A large
part of the $QK^\top$ time is spent writing the 64\,MiB matrix $S$ for the first time: it is
freshly allocated on every call and causes about 16{,}500 page faults per call, one per 4\,KiB
page, and in a test with $S$ already touched, $QK^\top$ dropped from 65\,ms to 40\,ms. Both of
these costs come from materialising $S$, which is what fusion in (c) removes.

\subsection*{(c) Fusion}

\texttt{attn\_fused.c} runs all three steps for one block of rows before moving to the next, so
only the scores of the current block are stored. I first tried literally one row at a time,
with an $N$-float score buffer. It passed, but at $d=128$ it took 726\,ms, slower than
\texttt{attn\_tiled.c}, because every row streamed all of $K^\top$ and $V$ (2\,MiB each, the
whole L3). Processing $B$ rows together reuses each pass over $K^\top$ and $V$ for $B$ rows, and
inside the block the steps use the same $8 \times 8$ register tiles as (b). I measured $B = 8$,
16, 32 and 64; $B=64$ was fastest or equal on every input (269\,ms at $d=128$, against 462\,ms
for $B=8$), so the score buffer is $64 \times N$ floats (1\,MiB at $N=4096$) instead of 64\,MiB.
Instead of dividing every score by the row sum, each output row is multiplied by $1/\text{sum}$
once at the end. This changes the rounding slightly, so \texttt{make check} now reports a small
non-zero error (at most $6.2\times10^{-6}$).

\begin{center}
\begin{tabular}{@{}lrrrrr@{}}
\toprule
Input & Time (ms) & GFLOP/s & \% of $P_1$ & vs.\ naive & vs.\ tiled \\
\midrule
$N=1024,\ d=64$  & 8.6   & 31.18 & 22.7 & $\times15.04$ & $\times1.33$ \\
$N=2048,\ d=64$  & 35.8  & 29.97 & 21.8 & $\times14.65$ & $\times1.17$ \\
$N=4096,\ d=64$  & 150.9 & 28.46 & 20.7 & $\times14.33$ & $\times1.50$ \\
$N=4096,\ d=128$ & 268.6 & 31.98 & 23.2 & $\times19.16$ & $\times1.24$ \\
\bottomrule
\end{tabular}
\end{center}

Removing the $N \times N$ matrix removes the costs that came with it: page faults at
$N=4096,\ d=64$ fall from 99{,}157 to 3{,}283 per 6 calls, and $QK^\top$ drops from 72\,ms to
48\,ms. Softmax also gets faster (70\,ms to 49\,ms), since its block of scores is now in cache
instead of DRAM. The three steps now take similar time (48\,ms $QK^\top$, 49\,ms softmax and
44\,ms $PV$), so the scalar \texttt{expf} in softmax is as large a cost as either matmul, and
the matmuls still run below half of peak.

\subsection*{(d) Best}

\texttt{attn\_best.c} starts from \texttt{attn\_fused.c} and makes two changes. The first is a
vectorised softmax. \texttt{expf} is a library call that clang cannot vectorise, so I wrote
\texttt{fast\_exp}, which computes $e^x = 2^n e^r$ with $n = \lfloor x\log_2 e + 0.5\rfloor$
and $|r| \le \ln 2/2$, evaluates $e^r$ with a degree-6 polynomial (Cephes \texttt{expf}
coefficients), and builds $2^n$ directly from the exponent bits. The row max and row sum are
kept in 8 independent lanes, so that neither is a single serial reduction. With this alone, the
softmax at $N=4096,\ d=64$ dropped from 49\,ms to 26\,ms and the whole kernel from 151\,ms to
about 114\,ms.

The second change is the online softmax used by FlashAttention (Milakov and Gimelshein, 2018).
Each block of 64 query rows now walks through the keys in tiles of $B_c$ columns, keeping a
running maximum $m$, sum $\ell$ and output $o$ per row. For each tile, the new maximum
$m'$ gives $\alpha = e^{m - m'}$, then $\ell \leftarrow \alpha\ell + \sum_j e^{s_j - m'}$ and
$o \leftarrow \alpha o + \sum_j e^{s_j - m'} v_j$, and $O = o/\ell$ after the last tile. The
score buffer shrinks from $64 \times N$ (1\,MiB at $N=4096$, larger than L2) to
$64 \times B_c$. With $B_c = 256$ the score tile and the matching slices of $K^\top$ and $V$
take 64\,KiB each and stay in the 512\,KiB L2, and L2 misses at $N=4096,\ d=64$ fell from
$108\times10^6$ to $3.8\times10^6$. Tiles of 128 and 512 columns were within a few percent of
256 and 1024 was 10--20\% slower. Row blocks of 16 to 128 changed $d=64$ by under 3\%; at
$d=128$, 16 and 32 rows were 9--30\% slower and 128 rows 3\% faster than 64, so I kept 64. \texttt{make check}
passes, with errors of at most $4.4\times10^{-6}$.

\begin{center}
\begin{tabular}{@{}lrrrrr@{}}
\toprule
Input & Time (ms) & GFLOP/s & \% of $P_1$ & vs.\ naive & vs.\ fused \\
\midrule
$N=1024,\ d=64$  & 5.9   & 45.33 & 32.9 & $\times21.87$ & $\times1.45$ \\
$N=2048,\ d=64$  & 23.3  & 46.13 & 33.5 & $\times22.54$ & $\times1.54$ \\
$N=4096,\ d=64$  & 95.7  & 44.86 & 32.6 & $\times22.59$ & $\times1.58$ \\
$N=4096,\ d=128$ & 180.4 & 47.62 & 34.6 & $\times28.53$ & $\times1.49$ \\
\bottomrule
\end{tabular}
\end{center}

What limits \texttt{attn\_best.c} now is the matmul code itself, at about a third of peak. The
$8 \times 8$ register tile written in plain C leaves the compiler to schedule loads and FMAs,
and changing its shape ($4 \times 16$, $2 \times 32$, $8 \times 16$) moved the time by under
5\%, while $16 \times 8$ needs more than the 16 vector registers and was $2.7\times$ slower.
Going further would need a hand-written micro-kernel with intrinsics.

% ===========================================================================
\section*{Q4: OpenMP}

\texttt{attn\_omp.c} is \texttt{attn\_best.c} with two OpenMP directives. The transpose of $K$
is split across threads with \texttt{parallel for}, and the loop over 64-row blocks is shared
with \texttt{omp for}, each thread allocating its own score tile. Each block writes its own rows
of $O$, so no synchronisation is needed, and the output is identical to \texttt{attn\_best.c}.
Times are medians of three runs after a warm-up run.

\begin{center}
\begin{tabular}{@{}r rrr rrr@{}}
\toprule
 & \multicolumn{3}{c}{$N=4096,\ d=64$} & \multicolumn{3}{c}{$N=4096,\ d=128$} \\
\cmidrule(lr){2-4}\cmidrule(lr){5-7}
Threads & Time (ms) & Speedup & Efficiency & Time (ms) & Speedup & Efficiency \\
\midrule
1 & 98.8 & 1.00 & 100\% & 184.8 & 1.00 & 100\% \\
2 & 50.1 & 1.97 & 99\%  & 91.4  & 2.02 & 101\% \\
3 & 35.7 & 2.77 & 92\%  & 73.3  & 2.52 & 84\%  \\
4 & 27.3 & 3.61 & 90\%  & 54.3  & 3.40 & 85\%  \\
6 & 27.7 & 3.57 & 60\%  & 52.4  & 3.53 & 59\%  \\
8 & 26.2 & 3.77 & 47\%  & 50.4  & 3.67 & 46\%  \\
\bottomrule
\end{tabular}
\end{center}

\noindent At $N=1024$ and $N=2048$ the speedups are similar ($\times3.81$ and $\times3.78$ on 4
threads, $\times3.90$ and $\times3.77$ on 8; all rows are in \texttt{results.csv}). With 8
threads the kernel reaches 164--177\,GFLOP/s, about a third of the all-core peak of
521\,GFLOP/s.

Scaling is close to linear up to 4 threads, one per physical core. It stays below $\times4$
partly because the clock falls from 4.24\,GHz with one busy core to 4.03--4.16\,GHz with four.
Three threads reach only 84--92\%, and pinning them to separate cores did not change this
(72.4\,ms against 72.6\,ms at $d=128$), so I cannot attribute it to thread placement. Beyond 4
threads there is almost no gain, because the extra threads are SMT siblings that share a core's
FMA units and load ports, and the kernel is limited by what each core can execute rather than
by memory. Placing 4
threads on only two physical cores (\texttt{OMP\_PLACES=threads}, \texttt{OMP\_PROC\_BIND=close})
makes $N=4096,\ d=64$ take 46\,ms instead of 25\,ms. The loop uses
\texttt{schedule(dynamic)}: with a static schedule, 6 threads gave every thread the same number
of blocks, so the threads sharing a core finished last and the run took 32\,ms, against 25\,ms
with a dynamic schedule. At 4 and 8 threads the two schedules performed the same.

\section*{Q5: Roofline}
\todo{}

\end{document}
